\section*{Appliquer une procédure}

\begin{exercise}[points=2]
    Sachant que les triangles $ABC$ et $A'B'C'$ sont semblables, complète le tableau ci-dessous.\\
    Cette question doit être réalisée sur la feuille d'examen.
    \begin{center}
        $\begin{tblr}{
                    hlines,
                    vlines,
                    columns = {1cm, c},
                    rows={1cm},
                    cell{1}{1}={r=1,c=3}{c},
                    cell{1}{4}={r=2,c=1}{c},
                    cell{1}{5}={r=1,c=3}{c}}
                ABC  &      &      & k & A'B'C' &        &        \\4
                |AB| & |BC| & |AC| &   & |A'B'| & |B'C'| & |A'C'| \\
                     & 12    & 16    &   & 9    &        & 24     \\
                20   & 36   &      &   &        & 9     & 10    \\
            \end{tblr}$
    \end{center}
\end{exercise}

\newpage
% 1 point par ligne correcte, 0,5 pt si une seule erreur par ligne
\begin{exercise}[points=2]
    Complète le tableau ci-dessous si tu sais que les solides $S$ et $S'$ sont semblables.\\
    Cette question doit être réalisée sur la feuille d'examen.\\
    P : périmète, A : aire, V : volume.
    \begin{center}
        $\begin{tblr}{
                    hlines,
                    vlines,
                    columns = {1cm, c},
                    rows={1cm},
                    cell{1}{1}={r=1,c=3}{c},
                    cell{1}{4}={r=2,c=1}{c},
                    cell{1}{5}={r=1,c=3}{c}}
                S  &    &   & k & S' &     &     \\
                P  & A  & V &   & P' 4& A'  & V'  \\
                40 & 80 &   &   & 20 &   & 8 \\
                10 &    &   & 3 &    & 90 & 270 \\
                %    & 6  &   &   &    & 54  &     \\
            \end{tblr}$
    \end{center}
\end{exercise}

% 1 point la proportion, 1 point pour la résolution
% TODO: changer les figures
\begin{exercise}[points=8]
    Détermine la valeur de $x$ en mobilisant la théorie adéquate (théorème de Thalès, triangle semblables ou triangles isométriques). Pour chaque exercice ton raisonnemet doit être complet et rigoureux.
    \begin{tasks}(2)
        %\task \begin{tikzpicture}[baseline=(current bounding box.north), scale=.8]
    \tikzmath{
        \angle = 30;
        \a = 1;
        \b = 2;
        \k = 3;
    }
    \tkzDefPoint(0,0){A}
    \tkzDefPoint(\a,0){B}
    \tkzDefPoint(\angle:\b){C}
    \tkzDefPointBy[homothety=center A ratio \k](B)
    \tkzGetPoint{D}
    \tkzDefLine[parallel= through D](B,C)
    \tkzInterLL(A,C)(D,tkzPointResult)
    \tkzGetPoint{E}

    \tkzDrawLines(A,E A,D)
    \tkzDrawLines[add=0 and 0](B,C D,E)
    \tkzDrawPoints(A,B,C,D,E)
    %labels
    \tkzLabelPoints[below](A)
    \tkzLabelPoints[below](B)
    \tkzLabelPoints[above](C)
    \tkzLabelPoints[below](D)
    \tkzLabelPoints[above](E)
    %valeurs
    \tkzLabelSegment[below](A,B){2}
    \tkzLabelSegment[below](B,D){5}
    \tkzLabelSegment[above](A,C){4}
    \draw[dashed] (A) to[bend left] (E);
    \tkzLabelSegment[above=30pt](A,E){$x$}
    %
    \draw (current bounding box.north)node{$CB\sslash ED$};

\end{tikzpicture}
        \task \input{tikz/thales02-v2.tikz}
        \task \begin{tikzpicture}[baseline=(current bounding box.north), scale=.8]
    \tikzmath{
        \angle = 40;
        \a = 2;
        \b = 2.5;
        \k = 2;
    }
    \tkzDefPoint(0,0){A}
    \tkzDefPoint(\a,0){B}
    \tkzDefPoint(\angle:\b){C}
    \tkzDefPointBy[homothety=center A ratio \k](B)
    \tkzGetPoint{D}
    \tkzDefLine[parallel= through D](B,C)
    \tkzInterLL(A,C)(D,tkzPointResult)
    \tkzGetPoint{E}

    \tkzDrawLines[add=0pt and 10 pt](A,E A,D)
    \tkzDrawLines[add=0 and 0](B,C D,E)

    \tkzMarkSegments[mark=||](A,B B,D)
    \tkzDrawPoints(A,B,C,D,E)
    %labels
    \tkzLabelPoints[below](A)
    \tkzLabelPoints[below](B)
    \tkzLabelPoints[above](C)
    \tkzLabelPoints[below](D)
    \tkzLabelPoints[above](E)
    %valeurs
    \tkzLabelSegment[right](B,C){x}
    \tkzLabelSegment[right](D,E){14}
    \tkzLabelSegment[below=5pt](A,B){6}

    \draw (current bounding box.north)node{$CB\sslash ED$};

\end{tikzpicture}
        %\task  \begin{tikzpicture}[baseline=(current bounding box.north), scale=0.8]
    \tkzDefPoints{0/0/A,3/0/B, 4/1/V}
    \tkzInterCC[R](A,2 cm)(B,4 cm)\tkzGetPoints{C}{Y}
    \tkzDrawPolygon(A,B,C)
    \tkzDefPointsBy[homothety=center A ratio 2.1](A,B,C){D,E,F}
    \tkzDefPointsBy[rotation=center A angle -60](D,E,F){D,E,F}
    \tkzDefPointsBy[translation=from A to V](D,E,F){D,E,F}
    \tkzDrawPolygon(D,E,F)
    \tkzDrawPoints(A,B,C,D,E,F)
    \tkzLabelPoints[below](A,B,E)
    \tkzLabelPoints[above](F,C)
    \tkzLabelPoints[left](D)
    \tkzLabelSegment[below](A,B){3}
    \tkzLabelSegment[above](B,C){4}
    \tkzLabelSegment[left](A,C){2}
    \tkzLabelSegment[below left](D,E){6,3}
    \tkzLabelSegment[right](E,F){8,4}
    \tkzLabelSegment[above left](D,F){4,2}
    \tkzLabelAngle[pos=0.6](B,A,C){\scriptsize 104\degres}
    \tkzLabelAngle[pos=0.6](E,D,F){$x$}
    \tkzMarkAngle[mark=](E,D,F)
\end{tikzpicture}
        \task \begin{tikzpicture}[baseline=(current bounding box.north),scale=1]
    \tkzDefPoints{0/0/A,4/0/B, 0/-1/V}
    \tkzInterCC[R](A,3 cm)(B,3 cm)\tkzGetPoints{C}{Y}
    \tkzDrawPolygon(A,B,C)
    \tkzDefPointsBy[homothety=center A ratio 0.6](A,B,C){D,E,F}
    \tkzDefPointsBy[rotation=center A angle 180](D,E,F){D,E,F}
    \tkzDefPointsBy[translation=from A to V](D,E,F){D,E,F}
    \tkzDrawPolygon(D,E,F)
    \tkzMarkSegments(A,C B,C)
    \tkzMarkSegments[mark=||](E,F F,D)
    \tkzDrawPoints(A,B,C,D,E,F)
    \tkzLabelPoints[below](A,B,F)
    \tkzLabelPoints[above](C,E,D)

    \tkzLabelSegment[below](A,B){5}
    \tkzLabelSegment[above left](A,C){3}
    \tkzLabelAngle[pos=0.7, sloped](B,A,C){\scriptsize 33,5\degres}

    \tkzLabelSegment[above](E,D){$x$}
    \tkzLabelSegment[right=1ex](D,F){1,8}
    \tkzLabelAngle[pos=0.5](D,F,E){\scriptsize 113\degres}
\end{tikzpicture}
        \task \begin{tikzpicture}[baseline=(current bounding box.north), rotate=30, scale=0.4]

    \tkzDefPoint(0,0){E}
    \tkzDefPoint(6,0){B}
    \tkzDefPoint(-6,0){A}
    \tkzDefPoint(0,-4){F}
    \tkzDefPoint(6,-4){H}
    \tkzDefPoint(-6,-4){G}
    \tkzDefPoint(-3,-4){C}
    \tkzDefPoint(3,-4){D}
    \tkzDrawLine(A,B)
    \tkzDrawLine(G,H)
    \tkzDrawSegment(E,C)
    \tkzDrawSegment(E,D)
    \tkzDrawSegment(A,C)
    \tkzDrawSegment(B,D)



    %\tkzDrawPolygon(A,B,C)
    \tkzLabelSegment[below](A,C){$x$}
    \tkzLabelSegment[right](B,D){6}
    %\tkzMarkAngle[mark=, size=0.5](B,A,C)
    %\tkzLabelAngle(B,A,C){70\degres}
    \tkzDrawPoints(A,B,C,D,E)
    \tkzLabelPoints[above](A)
    \tkzLabelPoints[above](B)
    %\tkzLabelPoints[above](C)
    \tkzLabelPoints[above](E)
    \tkzLabelPoints[below](C)
    \tkzLabelPoints[below](D)

    \tkzMarkSegments[mark=|](C,E E,D C,D)

    \draw (current bounding box.north)node[above]{$AB\sslash CD$ et $E$ est le milieu de $[AB]$};

    \end{tikzpicture}
    \end{tasks}
\end{exercise}

% TODO: mettre une étagère
\begin{exercise}[points=2]
    J'ai fixé deux planches au mur pour créer une étagère. Les planches sont-elles parallèles ? Note ton raisonnement.
    \begin{center}
        \begin{tikzpicture}[scale=2]
    \tkzDefPoints{-.5/0/A,0/5/B,0/1/C,2/1/D, 0/4/E,3/4/F}
    \draw (A) rectangle (B);
    \tkzDrawSegments(C,D E,F)
    \tkzDrawSegments[dashed](C,F D,E)
    \tkzInterLL(C,F)(D,E)
    \tkzGetPoint{G}
    \tkzDrawPoints(C,D,E,F,G)
    \tkzLabelSegment[fill=white,anchor=center, sloped](C,G){\SI{25}{cm}}
    \tkzLabelSegment[fill=white,anchor=center, sloped](D,G){\SI{20}{cm}}
    \tkzLabelSegment[fill=white,anchor=center, sloped](E,G){\SI{30}{cm}}
    \tkzLabelSegment[fill=white,anchor=center, sloped](F,G){\SI{35}{cm}}
\end{tikzpicture}
    \end{center}
\end{exercise}

